\ifdefined\gkver\else\def\gkver{student}\fi
\documentclass[\gkver]{gaokaozhenti}
\begin{document}

\chapter{2016年上海}

\begin{enumerate}
\item
%% number: 1
%% paperName: 2016年上海高考第 1 题 2 分
%% typeId: 1
%% type: 单选题
%% chapter: 原子和原子核
%% point: 原子核式结构
%% method: 无
%% score: 2
%% degree: 950
%% duplicateId: 0
%% body:
卢瑟福通过对 $\alpha$ 粒子散射实验结果的分析，提出了原子内部存在\xzanswer{D}
\fourchoices[answer=D]
{电子}
{中子}
{质子}
{原子核}
%% answer: D
%% memo:
\memoanswer{
卢瑟福由 $\alpha$ 粒子散射实验提出了原子的核式结构模型，故 D 正确。
}

\item
%% number: 2
%% paperName: 2016年上海高考第 2 题 2 分
%% typeId: 1
%% type: 单选题
%% chapter: 光学
%% point: 光的电磁说
%% method: 无
%% score: 2
%% degree: 850
%% duplicateId: 0
%% body:
一束单色光由空气进入水中，则该光在空气和水中传播时\xzanswer{D}
\fourchoices[answer=D]
{速度相同，波长相同}
{速度不同，波长相同}
{速度相同，频率相同}
{速度不同，频率相同}
%% answer: D
%% memo:
\memoanswer{
频率是光的特性，一束单色光由空气进入水中，频率不变，由 $v=\frac{c}{n}$ 知波速减小，由 $\lambda=\frac{v}{f}$ 知波长变短，故 A、B、C 错误，D 正确。
}

\item
%% number: 3
%% paperName: 2016年上海高考第 3 题 2 分
%% typeId: 1
%% type: 单选题
%% chapter: 光学
%% point: 光的电磁说
%% method: 无
%% score: 2
%% degree: 850
%% duplicateId: 0
%% body:
各种不同频率范围的电磁波按频率由大到小的排列顺序是\xzanswer{A}
\fourchoices[answer=A]
{$\gamma$ 射线、紫外线、可见光、红外线}
{$\gamma$ 射线、红外线、紫外线、可见光}
{紫外线、可见光、红外线、$\gamma$ 射线}
{红外线、可见光、紫外线、$\gamma$ 射线}
%% answer: A
%% memo:
\memoanswer{
电磁波谱中，按频率由高到低依次为 $\gamma$ 射线、紫外线、可见光、红外线，故 A 正确。
}

\item
%% number: 4
%% paperName: 2016年上海高考第 4 题 2 分
%% typeId: 1
%% type: 单选题
%% chapter: 运动和力
%% point: 牛顿第二定律
%% method: 无
%% score: 2
%% degree: 850
%% duplicateId: 0
%% body:
如图，顶端固定着小球的直杆固定在小车上，当小车向右做匀加速运动时，球所受合外力的方向沿图中的\xzanswer{D}
\onepicture[h=3.0cm]{\includesvg{figs/04}}
\fourchoices[answer=D]
{OA 方向}
{OB 方向}
{OC 方向}
{OD 方向}
%% answer: D
%% memo:
\memoanswer{
由牛顿第二定律，加速度的方向一定沿合外力的方向，所以合外力的方向一定沿 OD 方向，故 D 正确。
}

\item
%% number: 5
%% paperName: 2016年上海高考第 5 题 2 分
%% typeId: 1
%% type: 单选题
%% chapter: 电磁感应
%% point: 楞次定律
%% method: 无
%% score: 2
%% degree: 650
%% duplicateId: 0
%% body:
磁铁在线圈中心上方开始运动时，线圈中产生如图方向的感应电流，则磁铁\xzanswer{B}
\onepicture[h=3.4cm]{\includesvg{figs/05}}
\fourchoices[answer=B]
{向上运动}
{向下运动}
{向左运动}
{向右运动}
%% answer: B
%% memo:
\memoanswer{
由安培定则可知感应电流的磁场方向为竖直向上，故原磁场方向竖直向下，由楞次定律知，原磁场在增强，故磁铁向下运动，故 B 正确。
}

\item
%% number: 6
%% paperName: 2016年上海高考第 6 题 2 分
%% typeId: 1
%% type: 单选题
%% chapter: 原子和原子核
%% point: 天然放射性
%% method: 无
%% score: 2
%% degree: 650
%% duplicateId: 0
%% body:
放射性元素 A 经过 2 次 $\alpha$ 衰变和 1 次 $\beta$ 衰变后生成一新元素 B，则元素 B 在元素周期表中的位置较元素 A 的位置向前移动了\xzanswer{C}
\fourchoices[answer=C]
{1 位}
{2 位}
{3 位}
{4 位}
%% answer: C
%% memo:
\memoanswer{
A 元素经过两次 $\alpha$ 衰变，质子数减少 4，1 次 $\beta$ 衰变，质子数增加 1，故 B 元素在元素周期表中的位置较 A 元素向前移动了 3 位，故 C 正确。
}

\item
%% number: 7
%% paperName: 2016年上海高考第 7 题 2 分
%% typeId: 1
%% type: 单选题
%% chapter: 运动和力
%% point: 超重失重
%% method: 无
%% score: 2
%% degree: 550
%% duplicateId: 0
%% body:
在今年上海的某活动中引入了全国首个户外风洞飞行体验装置，体验者在风力作用下漂浮在半空。若减小风力，体验者在加速下落过程中\xzanswer{B}
\fourchoices[answer=B]
{失重且机械能增加}
{失重且机械能减少}
{超重且机械能增加}
{超重且机械能减少}
%% answer: B
%% memo:
\memoanswer{
当风力减小时，体验者加速下降，有向下的加速度，处于失重状态，风力做负功，体验者机械能减小，故 B 正确。
}

\item
%% number: 8
%% paperName: 2016年上海高考第 8 题 2 分
%% typeId: 1
%% type: 单选题
%% chapter: 磁场
%% point: 电流的磁场
%% method: 无
%% score: 2
%% degree: 550
%% duplicateId: 0
%% body:
如图，一束电子沿 $z$ 轴正向流动，则在图中 $y$ 轴上 A 点的磁场方向是\xzanswer{A}
\onepicture[h=3.0cm]{\includesvg{figs/08}}
\fourchoices[answer=A]
{$+ x$ 方向}
{$- x$ 方向}
{$+ y$ 方向}
{$- y$ 方向}
%% answer: A
%% memo:
\memoanswer{
电子向 $z$ 轴正向流动时，电流方向为 $z$ 轴负向，由安培定则，A 点的磁场方向为 $+x$ 方向，故 A 正确。
}

\item
%% number: 9
%% paperName: 2016年上海高考第 9 题 3 分
%% typeId: 1
%% type: 单选题
%% chapter: 光学
%% point: 光的干涉
%% method: 无
%% score: 3
%% degree: 650
%% duplicateId: 0
%% body:
在双缝干涉实验中，屏上出现了明暗相间的条纹，则\xzanswer{D}
\fourchoices[answer=D]
{中间条纹间距较两侧更宽}
{不同色光形成的条纹完全重合}
{双缝间距离越大条纹间距离也越大}
{遮住一条缝后屏上仍有明暗相间的条纹}
%% answer: D
%% memo:
\memoanswer{
双缝干涉实验中条纹间距相等，故 A 错误；由 $\Delta x = \frac{L}{d}\lambda$，不同色光波长不同，条纹间距不同，故不同色光形成的条纹不重合，故 B 错误；由 $\Delta x = \frac{L}{d}\lambda$，双缝间距 $d$ 越小条纹间距越大，故 C 错误；遮住一条缝后，双缝干涉条纹变为单缝衍射条纹，故 D 正确。
}

\item
%% number: 10
%% paperName: 2016年上海高考第 10 题 3 分
%% typeId: 1
%% type: 单选题
%% chapter: 原子和原子核
%% point: 天然放射性
%% method: 无
%% score: 3
%% degree: 550
%% duplicateId: 0
%% body:
研究放射性元素射线性质的实验装置如图所示。两块平行放置的金属板 A、B 分别与电源的两极 a、b 连接，放射源发出的射线从其上方小孔向外射出。则\xzanswer{B}
\onepicture[h=3.6cm]{\includesvg{figs/10}}
\fourchoices[answer=B]
{a 为电源正极，到达 A 板的为 $\alpha$ 射线}
{a 为电源正极，到达 A 板的为 $\beta$ 射线}
{a 为电源负极，到达 A 板的为 $\alpha$ 射线}
{a 为电源负极，到达 A 板的为 $\beta$ 射线}
%% answer: B
%% memo:
\memoanswer{
$\alpha$ 射线速度为 $\beta$ 射线速度的 $\frac{1}{10}$，由两者在极板间偏转距离相同，则 $s = v_{0}t = v_{0}\sqrt{\frac{md}{2qE}}$，因 $\alpha$ 粒子的比荷是 $\beta$ 粒子的近两千分之一，故偏向 B 极板的射线为 $\alpha$ 射线，a 为电源正极，电场线方向向右，到达 A 极板的射线为 $\beta$ 射线，故 B 正确。
}

\item
%% number: 11
%% paperName: 2016年上海高考第 11 题 3 分
%% typeId: 1
%% type: 单选题
%% chapter: 运动和力
%% point: 力学单位制
%% method: 无
%% score: 3
%% degree: 550
%% duplicateId: 0
%% body:
国际单位制中，\textbf{不是}电场强度的单位是\xzanswer{C}
\fourchoices[answer=C]
{$\UNC$}
{$\UVm$}
{$\UJC$}
{$\UT\cdot\Ums$}
%% answer: C
%% memo:
\memoanswer{
由 $E=\frac{F}{q}$，$E=\frac{U}{d}$，故 A、B 正确；由 $qE=qvB$ 得 $E=vB$，故 D 正确；故不是电场强度单位的是选项 C。
}

\item
%% number: 12
%% paperName: 2016年上海高考第 12 题 3 分
%% typeId: 1
%% type: 单选题
%% chapter: 气体性质
%% point: 液柱移动定性分析
%% method: 无
%% score: 3
%% degree: 450
%% duplicateId: 0
%% body:
如图，粗细均匀的玻璃管 A 和 B 由一橡皮管连接，一定质量的空气被水银柱封闭在 A 管内，初始时两管水银面等高，B 管上方与大气相通。若固定 A 管，将 B 管沿竖直方向缓慢下移一小段距离 $H$，A 管内的水银面高度相应变化 $h$，则\xzanswer{B}
\onepicture[h=3.6cm]{\includesvg{figs/12}}
\fourchoices[answer=B]
{$h = H$}
{$h < \frac{H}{2}$}
{$h = \frac{H}{2}$}
{$\frac{H}{2} < h < H$}
%% answer: B
%% memo:
\memoanswer{
初始时 A 管中气体压强为 $p_{0}$，B 管沿竖直方向下移一小段距离 $H$，假设气体压强与体积无关，则 A 管内水银面下降 $\frac{H}{2}$，实际上气体压强随体积增大而减小，即 $p < p_{0}$，故水银面高度变化 $h < \frac{H}{2}$，故 B 正确。
}

\item
%% number: 13
%% paperName: 2016年上海高考第 13 题 3 分
%% typeId: 1
%% type: 单选题
%% chapter: 电路
%% point: 电动势
%% method: 无
%% score: 3
%% degree: 550
%% duplicateId: 0
%% body:
电源电动势反映了电源把其它形式的能量转化为电能的能力，因此\xzanswer{C}
\fourchoices[answer=C]
{电动势是一种非静电力}
{电动势越大，表明电源储存的电能越多}
{电动势的大小是非静电力做功能力的反映}
{电动势就是闭合电路中电源两端的电压}
%% answer: C
%% memo:
\memoanswer{
电源内部为非静电力做功，但电动势不是非静电力，故 A 错误；电动势越大，表明在电源内部移动等量的电荷时，非静电力做功多，故 B 错误，C 正确；闭合电路中，因电源存在内阻，电动势大于路端电压，故 D 错误。
}

\item
%% number: 14
%% paperName: 2016年上海高考第 14 题 3 分
%% typeId: 1
%% type: 单选题
%% chapter: 直线运动
%% point: 匀变速直线运动
%% method: 无
%% score: 3
%% degree: 450
%% duplicateId: 0
%% body:
物体做匀加速直线运动，相继经过两段距离均为 $16\Um$ 的路程，第一段用时 $4\Us$，第二段用时 $2\Us$，则物体的加速度是\xzanswer{B}
\fourchoices[answer=B]
{$\frac{2}{3}\Umsq$}
{$\frac{4}{3}\Umsq$}
{$\frac{8}{9}\Umsq$}
{$\frac{16}{9}\Umsq$}
%% answer: B
%% memo:
\memoanswer{
由中间时刻的瞬时速度等于这段时间的平均速度。第一段距离中间时刻的速度 $v_{2}=4\Ums$，第二段距离中间时刻的速度 $v_{5}=8\Ums$。由 $a=\frac{v_{5}-v_{2}}{t}=\frac{4}{3}\Umsq$，故 B 正确。
}

\item
%% number: 15
%% paperName: 2016年上海高考第 15 题 3 分
%% typeId: 1
%% type: 单选题
%% chapter: 力物体的平衡
%% point: 力矩平衡
%% method: 无
%% score: 3
%% degree: 450
%% duplicateId: 0
%% body:
如图，始终竖直向上的力 $F$ 作用在三角板 A 端，使其绕 B 点在竖直平面内缓慢地沿顺时针方向转动一小角度，力 $F$ 对 B 点的力矩为 $M$，则转动过程中\xzanswer{A}
\onepicture[h=2.8cm]{\includesvg{figs/15}}
\fourchoices[answer=A]
{$M$ 减小，$F$ 增大}
{$M$ 减小，$F$ 减小}
{$M$ 增大，$F$ 增大}
{$M$ 增大，$F$ 减小}
%% answer: A
%% memo:
\memoanswer{
拉力力矩始终和重力力矩相等，由于重力力矩减小，故拉力力矩减小，因重力臂小于拉力臂，重力不变，故拉力 $F$ 增大，故 A 正确。
}

\item
%% number: 16
%% paperName: 2016年上海高考第 16 题 3 分
%% typeId: 1
%% type: 单选题
%% chapter: 曲线运动
%% point: 圆周运动的应用
%% method: 无
%% score: 3
%% degree: 350
%% duplicateId: 0
%% body:
风速仪结构如图~\ref{2016上海16a} 所示。光源发出的光经光纤传输，被探测器接收，当风轮旋转时，通过齿轮带动凸轮圆盘旋转，当圆盘上的凸轮经过透镜系统时光被遮挡。已知风轮叶片转动半径为 $r$，每转动 $n$ 圈带动凸轮圆盘转动一圈。若某段时间 $\Delta t$ 内探测器接收到的光强随时间变化关系如图~\ref{2016上海16b} 所示，则该时间段内风轮叶片\xzanswer{B}
\twopicture[labelA=2016上海16a, labelB=2016上海16b, numA=（a）, numB=（b）, wA=4.6cm, wB=5.0cm]
{\includesvg{figs/16a}}
{\includesvg{figs/16b}}
\fourchoices[answer=B]
{转速逐渐减小，平均速率为 $\frac{4\pi nr}{\Delta t}$}
{转速逐渐减小，平均速率为 $\frac{8\pi nr}{\Delta t}$}
{转速逐渐增大，平均速率为 $\frac{4\pi nr}{\Delta t}$}
{转速逐渐增大，平均速率为 $\frac{8\pi nr}{\Delta t}$}
%% answer: B
%% memo:
\memoanswer{
挡光时间间隔越来越长，故风速仪转速逐渐减小；$\Delta t$ 时间内，光强为 4 个周期，风速仪转动的弧长为 $4n\cdot2\pi r$，故平均速率为 $\frac{8\pi nr}{\Delta t}$，故 B 正确。
}

\item
%% number: 17
%% paperName: 2016年上海高考第 17 题 4 分
%% typeId: 2
%% type: 多选题
%% chapter: 气体性质
%% point: 阿伏伽德罗常数
%% method: 无
%% score: 4
%% degree: 450
%% duplicateId: 0
%% body:
某气体的摩尔质量为 $M$，分子质量为 $m$。若 1 摩尔该气体的体积为 $V_{m}$，密度为 $\rho$，则该气体单位体积分子数为（阿伏伽德罗常数为 $N_{\mathrm{A}}$）\xzanswer{ABC}
\fourchoices[answer=ABC]
{$\frac{N_{\mathrm{A}}}{V_{m}}$}
{$\frac{M}{mV_{m}}$}
{$\frac{\rho N_{\mathrm{A}}}{M}$}
{$\frac{\rho N_{\mathrm{A}}}{m}$}
%% answer: ABC
%% memo:
\memoanswer{
单位体积内的分子数为 $\frac{N_{\mathrm{A}}}{V_{m}}$，故 A 正确；因 $N_{\mathrm{A}} = \frac{M}{m}$，故 $\frac{M}{mV_{m}}$ 也表示单位体积内的分子数，故 B 正确；因 $V_{m} = \frac{M}{\rho}$，故 $\frac{\rho N_{\mathrm{A}}}{M}$ 同样表示单位体积内的分子数，故 C 正确。
}

\item
%% number: 18
%% paperName: 2016年上海高考第 18 题 4 分
%% typeId: 2
%% type: 多选题
%% chapter: 电路
%% point: 动态分析
%% method: 无
%% score: 4
%% degree: 350
%% duplicateId: 0
%% body:
如图所示电路中，电源内阻忽略不计。闭合电键，电压表示数为 $U$，电流表示数为 $I$；在滑动变阻器 $R_{1}$ 的滑片 P 由 a 端滑到 b 端的过程中\xzanswer{BC}
\onepicture[h=4.6cm]{\includesvg{figs/18}}
\fourchoices[answer=BC]
{$U$ 先变大后变小}
{$I$ 先变小后变大}
{$U$ 与 $I$ 比值先变大后变小}
{$U$ 变化量与 $I$ 变化量比值等于 $R_{3}$}
%% answer: BC
%% memo:
\memoanswer{
电源内阻不计，电压表读数不变，故 A 错误；滑动变阻器滑片由 a 端滑到 b 端，电阻 $R_{1}$ 先增后减，电流表示数先减后增，故 B 正确；由 $\frac{U}{I} = R_{2} + R_{1}$，电阻 $R_{1}$ 先增后减，故 $U$ 与 $I$ 比值先增后减，故 C 正确；$\frac{\Delta U}{\Delta I}$ 与 $R_{3}$ 无关，故 D 错误。
}

\item
%% number: 19
%% paperName: 2016年上海高考第 19 题 4 分
%% typeId: 2
%% type: 多选题
%% chapter: 电磁感应
%% point: 二级电磁感应
%% method: 无
%% score: 4
%% degree: 350
%% duplicateId: 0
%% body:
如图~\ref{2016上海19a}，螺线管内有平行于轴线的外加匀强磁场，以图中箭头所示方向为其正方向。螺线管与导线框 abcd 相连，导线框内有一小金属圆环 L，圆环与导线框在同一平面内。当螺线管内的磁感应强度 $B$ 随时间按图~\ref{2016上海19b} 所示规律变化时\xzanswer{AD}
\twopicture[labelA=2016上海19a, labelB=2016上海19b, numA=（a）, numB=（b）, wA=4.2cm, wB=5.2cm]
{\includesvg{figs/19a}}
{\includesvg{figs/19b}}
\fourchoices[answer=AD]
{在 $t_{1} \sim t_{2}$ 时间内，L 有收缩趋势}
{在 $t_{2} \sim t_{3}$ 时间内，L 有扩张趋势}
{在 $t_{2} \sim t_{3}$ 时间内，L 内有逆时针方向的感应电流}
{在 $t_{3} \sim t_{4}$ 时间内，L 内有顺时针方向的感应电流}
%% answer: AD
%% memo:
\memoanswer{
在 $t_{1} \sim t_{2}$ 的时间内，$B$ 的斜率逐渐增大，回路内的感应电流逐渐增大，回路电流产生的磁场增大，由楞次定律的广义表述，L 有收缩的趋势，故 A 正确；在 $t_{2} \sim t_{3}$ 的时间内，$B$ 的斜率不变，回路内的感应电流不变，回路电流产生的磁场不变，L 中没有感应电流产生，故 L 既没有收缩的趋势，也没有扩张的趋势，故 B、C 错误；在 $t_{3} \sim t_{4}$ 的时间内，负向磁场 $B$ 的斜率减小，回路内的顺时针方向的感应电流减小，回路内感应电流向里的磁场减小，由楞次定律可知，L 内有顺时针方向的感应电流，故 D 正确。
}

\item
%% number: 20
%% paperName: 2016年上海高考第 20 题 4 分
%% typeId: 2
%% type: 多选题
%% chapter: 机械波
%% point: 波的叠加
%% method: 无
%% score: 4
%% degree: 350
%% duplicateId: 0
%% body:
甲、乙两列横波在同一介质中分别从波源 M、N 两点沿 $x$ 轴相向传播，波速为 $2\Ums$，振幅相同；某时刻的图像如图所示。则\xzanswer{ABD}
\onepicture[w=9.5cm]{\includesvg{figs/20}}
\fourchoices[answer=ABD]
{甲乙两波的起振方向相反}
{甲乙两波的频率之比为 $3:2$}
{再经过 $3\Us$，平衡位置在 $x=7\Um$ 处的质点振动方向向下}
{再经过 $3\Us$，两波源间（不含波源）有 5 个质点位移为零}
%% answer: ABD
%% memo:
\memoanswer{
甲、乙两列波起振方向一下一上，故 A 正确；

$\lambda_{\text{甲}}=4\Um$，$\lambda_{\text{乙}}=6\Um$，在同种介质中机械波波速相同，频率 $f=\frac{v}{\lambda}$，故 $\frac{f_{\text{甲}}}{f_{\text{乙}}}=\frac{3}{2}$，故 B 正确；

再经过 $3\Us$，甲、乙波都传播 $6\Um$，$x=7\Um$ 处，甲波在波谷处，乙波在 $x$ 轴上方向上振动但不在波峰处，故该质点振动方向向上，故 C 错误；

画出波形图，由波的叠加可知两波源间有 5 个质点位移为 0，故 D 正确。
}

\item
%% number: 21
%% paperName: 2016年上海高考第 21 题 4 分
%% typeId: 3
%% type: 填空题
%% chapter: 磁场
%% point: 磁感应强度
%% method: 无
%% score: 4
%% degree: 750
%% duplicateId: 0
%% body:
形象描述磁场分布的曲线叫做\tkanswer[2.0cm]{磁感线}，通常\tkanswer[2.4cm]{磁感应强度}的大小也叫做磁通量密度。
%% answer: 磁感线，磁感应强度
%% memo:
\memoanswer{
磁感线是形象化描述磁场的曲线，由 $B=\frac{\Phi}{S}$，磁感应强度 $B$ 又叫磁通量密度。
}

\item
%% number: 22
%% paperName: 2016年上海高考第 22 题 4 分
%% typeId: 3
%% type: 填空题
%% chapter: 运动和力
%% point: 动量守恒定律
%% method: 无
%% score: 4
%% degree: 650
%% duplicateId: 0
%% body:
如图，粗糙水平面上，两物体 A、B 以轻绳相连，在恒力 $F$ 作用下做匀速运动。某时刻轻绳断开，A 在 $F$ 牵引下继续前进，B 最后静止。则在 B 静止前，A 和 B 组成的系统动量\tkanswer[1.4cm]{守恒}（选填：“守恒”或“不守恒”）；在 B 静止后，A 和 B 组成的系统动量\tkanswer[1.4cm]{不守恒}（选填：“守恒”或“不守恒”）。
\onepicture[w=5.0cm, align=right]{\includesvg{figs/22}}
%% answer: 守恒，不守恒
%% memo:
\memoanswer{
B 在停止运动前，其所受滑动摩擦力不变，系统所受合外力为 0，动量守恒，B 停止运动后，其所受滑动摩擦力变为 0，系统所受合外力不为 0，动量不守恒。
}

\item
%% number: 22
%% paperName: 2016年上海高考第 22 题 4 分
%% typeId: 3
%% type: 填空题
%% chapter: 曲线运动
%% point: 人造地球卫星
%% method: 无
%% score: 4
%% degree: 550
%% duplicateId: 0
%% body:
两颗卫星绕地球运行的周期之比为 $27:1$，则它们的角速度之比为\tkanswer[1.6cm]{$1:27$}，轨道半径之比为\tkanswer[1.6cm]{$9:1$}。
%% answer: $1:27$，$9:1$
%% memo:
\memoanswer{
由 $T=\frac{2\pi}{\omega}$，角速度比为周期比的反比，故角速度比为 $1:27$，由开普勒第三定律 $T^{2}\propto r^{3}$，可得半径比为 $9:1$。
}

\item
%% number: 23
%% paperName: 2016年上海高考第 23 题 4 分
%% typeId: 3
%% type: 填空题
%% chapter: 曲线运动
%% point: 平抛运动
%% method: 无
%% score: 4
%% degree: 450
%% duplicateId: 0
%% body:
如图，圆弧形凹槽固定在水平地面上，其中 ABC 是位于竖直平面内以 O 为圆心的一段圆弧，OA 与竖直方向的夹角为 $\alpha$。一小球以速度 $v_{0}$ 从桌面边缘 P 水平抛出，恰好从 A 点沿圆弧的切线方向进入凹槽。小球从 P 到 A 的运动时间为\tkanswer[2.6cm]{$\frac{v_{0}\tan\alpha}{g}$}；直线 PA 与竖直方向间夹角 $\beta=$\tkanswer[3.2cm]{$\arctan(2\cot\alpha)$}。
\onepicture[w=5.5cm, align=right]{\includesvg{figs/23}}
%% answer: $\frac{v_{0}\tan\alpha}{g}$，$\arctan(2\cot\alpha)$
%% memo:
\memoanswer{
平抛运动速度偏向角为 $\alpha$，且 $\tan\alpha=\frac{gt}{v_{0}}$，解得 $t=\frac{v_{0}\tan\alpha}{g}$，由速度偏向角和位移偏向角的关系，得 $\tan\alpha=2\tan\left(\frac{\pi}{2}-\beta\right)$，解得 $\beta=\arctan(2\cot\alpha)$。
}

\item
%% number: 24
%% paperName: 2016年上海高考第 24 题 4 分
%% typeId: 3
%% type: 填空题
%% chapter: 电场
%% point: 带电粒子的圆周运动
%% method: 无
%% score: 4
%% degree: 350
%% duplicateId: 0
%% body:
如图，质量为 $m$ 的带电小球 A 用绝缘细线悬挂于 O 点，处于静止状态。施加一水平向右的匀强电场后，A 向右摆动，摆动的最大角度为 $60^\circ$，则 A 受到的电场力大小为\tkanswer[2.4cm]{$\frac{\sqrt{3}}{3}mg$}。在改变电场强度的大小和方向后，小球 A 的平衡位置在 $\alpha=60^\circ$ 处，然后再将 A 的质量改变为 $2m$，其新的平衡位置在 $\alpha=30^\circ$ 处，A 受到的电场力大小为\tkanswer[1.6cm]{$mg$}。
\onepicture[h=3.2cm, align=right]{\includesvg{figs/24}}
%% answer: $\frac{\sqrt{3}}{3}mg$，$mg$
%% memo:
\memoanswer{
由动能定理得 $-mgl(1-\cos60^\circ)+qE\sin60^\circ=0$，可得 $qE=\frac{\sqrt{3}mg}{3}$；由题意，改变小球电场强度大小和方向后，设电场方向与竖直方向的夹角为 $\beta$，小球受力图如图甲所示，由正弦定理 $\frac{qE}{\sin60^\circ}=\frac{mg}{\sin(180^\circ-60^\circ-\beta)}$，改变小球质量后，小球受力图如图乙所示，由正弦定理得 $\frac{qE}{\sin30^\circ}=\frac{2mg}{\sin(180^\circ-30^\circ-\beta)}$，解得 $\beta=60^\circ$，$qE=mg$。

\twopicture[numA=甲, numB=乙, hA=2.6cm, hB=3.0cm]
{figs/24_an1.png}
{figs/24_an2.png}
}

\item
%% number: 25
%% paperName: 2016年上海高考第 25 题 4 分
%% typeId: 3
%% type: 填空题
%% chapter: 功和能
%% point: 功能中的图象
%% method: 无
%% score: 4
%% degree: 350
%% duplicateId: 0
%% body:
地面上物体在变力 $F$ 作用下由静止开始竖直向上运动，力 $F$ 随高度 $x$ 的变化关系如图所示，物体能上升的最大高为 $h$，$h < H$。当物体加速度最大时其高度为\tkanswer[2.2cm]{0 或 $h$}，加速度的最大值为\tkanswer[2.6cm]{$\frac{gh}{2H-h}$}。
\onepicture[w=4.8cm, align=right]{\includesvg{figs/25}}
%% answer: 0 或 $h$，$\frac{gh}{2H-h}$
%% memo:
\memoanswer{
由力 $F$ 关于高度 $x$ 线性变化可知加速度 $a$ 也随高度 $x$ 线性变化，因为物体从静止开始运动且最后速度为 0，故在运动中间位置 $a=0$，且各点加速度大小关于此点对称，故最高点最低点加速度最大，物体所在高度为 0 或 $h$；

由图像可知外力 $F = F_{0} - \frac{F_{0}}{H}x$；在最低点 $F_{0} - mg = ma$，在最高点 $F = F_{0} - \frac{F_{0}}{H}h$，$mg - F = ma$，解方程得 $a = \frac{gh}{2H-h}$。
}

\item
%% number: 26
%% paperName: 2016年上海高考第 26 题 3 分
%% typeId: 4
%% type: 实验题
%% chapter: 功和能
%% point: DIS研究机械能守恒定律
%% method: 无
%% score: 3
%% degree: 750
%% duplicateId: 0
%% body:
在“用 DIS 研究机械能守恒定律”的实验中，用到的传感器是\tkanswer[2.2cm]{光电门}传感器。若摆锤直径的测量值大于其真实值会造成摆锤动能的测量值偏\tkanswer[1.2cm]{大}。（选填：“大”或“小”）。
%% answer: 光电门，大
%% memo:
\memoanswer{
用光电门传感器由公式 $v = \frac{s}{t}$ 测得摆锤的速度，摆锤直径的测量值偏大，则计算得到摆锤的速度偏大，摆锤动能测量值偏大。
}

\item
%% number: 27
%% paperName: 2016年上海高考第 27 题 6 分
%% typeId: 4
%% type: 实验题
%% chapter: 电路
%% point: 多用表的使用
%% method: 无
%% score: 6
%% degree: 550
%% duplicateId: 0
%% body:
在“用多用电表测电阻、电流和电压”的实验中
\begin{enumerate}
	\item
	（多选题）用多用电测电流或电阻的过程中
	\fourchoices[answer=AD]
	{在测量电阻时，更换倍率后必须重新进行调零}
	{在测量电流时，更换量程后必须重新进行调零}
	{在测量未知电阻时，必须先选择倍率最大挡进行试测}
	{在测量未知电流时，必须先选择电流最大量程进行试测}
	\item
	测量时多用电表指针指在如图所示位置。若选择开关处于“$10\UV$”挡，其读数为\tkanswer[1.6cm]{5.4}$\UV$；若选择开关处于“$\times10$”挡，其读数为\tkanswer[1.4cm]{小于}$200\UO$（选填：“大于”，“等于”或“小于”）。
\end{enumerate}
\onepicture[h=4.6cm, align=right]{figs/27.png}
%% answer: （1）AD（2）5.4，小于
\jdanswer{
\begin{enumerate}
	\item
	AD

	\item
	5.4，小于

\end{enumerate}
}
%% memo:
\memoanswer{
(1) 在测电阻时，更换倍率必须重新进行电阻调零，测电流则不需要，故 A 正确，B 错误；多用电表测电阻时，不管使用多大倍率，电阻最大值都为无穷大，不需用最大倍率测试，选择合适倍率尝试即可，测电流则需要最大量程试测，故 C 错误，D 正确。

(2) 若选择开关位于“$10\UV$”挡，则由读数规则可知其读数为“$5.4\UV$”；若选择开关处于“$\times10$”挡，由多用电表表盘上刻度分布规律可知，其读数小于 $200\UO$。
}

\item
%% number: 28
%% paperName: 2016年上海高考第 28 题 7 分
%% typeId: 4
%% type: 实验题
%% chapter: 电场
%% point: 描绘等势面
%% method: 无
%% score: 7
%% degree: 450
%% duplicateId: 0
%% body:
“用 DIS 描绘电场的等势线”的实验装置示意图如图所示。
\begin{enumerate}
	\item
	（单选题）该实验描绘的是
	\fourchoices[answer=B]
	{两个等量同种电荷周围的等势线}
	{两个等量异种电荷周围的等势线}
	{两个不等量同种电荷周围的等势线}
	{两个不等量异种电荷周围的等势线}
	\item
	（单选题）实验操作时，需在平整的木板上依次铺放
	\fourchoices[answer=D]
	{导电纸、复写纸、白纸}
	{白纸、导电纸、复写纸}
	{导电纸、白纸、复写纸}
	{白纸、复写纸、导电纸}
	\item
	若电压传感器的红、黑探针分别接触图中 d、f 两点（f、d 连线与 A、B 连线垂直）时，示数小于零。为使示数为零，应保持红色探针与 d 点接触，而将黑色探针\tkanswer[1.4cm]{向右}（选填：“向左”或“向右”）移动。
\end{enumerate}
\onepicture[w=4.8cm, align=right]{\includesvg{figs/28}}
%% answer: （1）B（2）D（3）向右
\jdanswer{
\begin{enumerate}
	\item
	B

	\item
	D

	\item
	向右

\end{enumerate}
}
%% memo:
\memoanswer{
(1) 描绘静电场等势线时，用直流电源来描绘等量异种点电荷的等势线，故 B 正确。

(2) 在平整的木板上，从下到上依次为白纸、复写纸、导电纸，故 D 正确。

(3) 电压传感器的红、黑探针依次接触 d、f，示数小于零，则 $\varphi_{d}<\varphi_{f}$，由等势面的特点，为使示数为 0，应将黑探针移至高电势处，即黑探针 f 右移。
}

\item
%% number: 29
%% paperName: 2016年上海高考第 29 题 8 分
%% typeId: 4
%% type: 实验题
%% chapter: 气体性质
%% point: 盖吕萨克定律
%% method: 无
%% score: 8
%% degree: 350
%% duplicateId: 0
%% body:
某同学制作了一个结构如图~\ref{2016上海29a} 所示的温度计。一端封闭的轻质细管可绕封闭端 O 自由转动，管长 $0.5\Um$。将一量程足够大的力传感器调零，细管的开口端通过细线挂于力传感器挂钩上，使细管保持水平、细线沿竖直方向。在气体温度为 $270\UK$ 时，用一段水银将长度为 $0.3\Um$ 的气柱封闭在管内。实验时改变气体温度，测得封闭气柱长度 $l$ 和力传感器读数 $F$ 之间的关系如图~\ref{2016上海29b} 所示（实验中大气压强不变）。
\begin{enumerate}
	\item
	管内水银柱长度为\tkanswer[1.4cm]{0.1}$\Um$，为保证水银不溢出，该温度计能测得的最高温度为\tkanswer[1.6cm]{360}$\UK$。
	\item
	若气柱初始长度大于 $0.3\Um$，该温度计能测量的最高温度将\tkanswer[1.4cm]{减小}（选填：“增大”，“不变”或“减小”）。
	\item
	若实验中大气压强略有升高，则用该温度计测出的温度将\tkanswer[1.4cm]{偏低}（选填：“偏高”，“不变”或“偏低”）。
\end{enumerate}
\twopicture[labelA=2016上海29a, labelB=2016上海29b, numA=（a）, numB=（b）, wA=4.6cm, wB=4.8cm, align=right]
{\includesvg{figs/29a}}
{\includesvg{figs/29b}}
%% answer: （1）0.1，360（2）减小（3）偏低
\jdanswer{
\begin{enumerate}
	\item
	0.1，360

	\item
	减小

	\item
	偏低

\end{enumerate}
}
%% memo:
\memoanswer{
(1) 设水银的长度为 $2b$，则由力矩平衡得

$Fl_{0}=mg(l+b)$，

故 $F=\frac{mg}{l_{0}}l+\frac{mgb}{l_{0}}$，

从图象得出 $F-l$ 的关系 $F=l+0.05$，

对比系数得 $\frac{mg}{l_{0}}=1$，$\frac{mgb}{l_{0}}=0.05$，

解得 $2b = 0.1\Um$，

由等压变化得

$\frac{l}{T_{0}}=\frac{l_{0}-2b}{T}$，

代入数据得 $T = 360\UK$。

(2) 若气柱初始长度大于 $0.3\Um$，则气体体积变化幅度变小，能测量的最高温度减小。

(3) 由 $\frac{V}{T}=\frac{\Delta V}{\Delta T}=\frac{k}{p}$，大气压强略有升高后，比值变小，升高相同的温度，气体体积变化小，示数偏低。
}

\item
%% number: 30
%% paperName: 2016年上海高考第 30 题 10 分
%% typeId: 5
%% type: 辨析题
%% chapter: 气体性质
%% point: 液柱移动定性分析
%% method: 无
%% score: 10
%% degree: 550
%% duplicateId: 0
%% body:
如图，两端封闭的直玻璃管竖直放置，一段水银将管内气体分隔为上下两部分 A 和 B，上下两部分气体初温度相等，且体积 $V_{A} > V_{B}$。
\begin{enumerate}
	\item
	若 A、B 两部分气体同时升高相同的温度，水银柱将如何移动？

	某同学解答如下：

	设两部分气体压强不变，由

	$\frac{V_{1}}{T_{1}}=\frac{V_{2}}{T_{2}}$，$\ldots$，$\Delta V=\frac{\Delta T}{T}V$，$\ldots$，所以水银柱将向下移动。

	上述解答是否正确？若正确，请写出完整的解答；若不正确，请说明理由并给出正确的解答。
	\item
	在上下两部分气体升高相同温度的过程中，水银柱位置发生变化，最后稳定在新的平衡位置，A、B 两部分气体始末状态压强的变化量分别为 $\Delta p_{A}$ 和 $\Delta p_{B}$，分析并比较二者的大小关系。
\end{enumerate}
\onepicture[h=4.6cm, align=right]{\includesvg{figs/30}}
%% answer: （1）不正确，水银柱向上移动（2）$\Delta p_{A}=\Delta p_{B}$
\jdanswer{
\begin{enumerate}
	\item
	不正确。水银柱移动的原因是升温后，由于压强变化造成受力平衡被破坏，因此应该假设气体体积不变，由压强变化判断移动方向；正确解法见详解，水银柱向上移动。

	\item
	$\Delta p_{A}=\Delta p_{B}$

\end{enumerate}
}
%% memo:
\memoanswer{
(1) 不正确。

水银柱移动的原因是升温后，由于压强变化造成受力平衡被破坏，因此应该假设气体体积不变，由压强变化判断移动方向。正确解法：设升温后上下部分气体体积不变，则由查理定律得

$\frac{p}{T}=\frac{p'}{T+\Delta T}$，

$\Delta p = p' - p = \frac{\Delta T}{T}p$，

因为 $\Delta T > 0$，$p_{A} < p_{B}$，可知 $\Delta p_{A} < \Delta p_{B}$，所以水银柱向上移动。

(2) 升温前有

$p_{B}=p_{A}+p_{h}$（$p_{h}$ 为汞柱压强），

升温后同样有

$p_{B}' = p_{A}' + p_{h}$，

两式相减可得 $\Delta p_{A} = \Delta p_{B}$。
}

\item
%% number: 31
%% paperName: 2016年上海高考第 31 题 12 分
%% typeId: 6
%% type: 计算题
%% chapter: 功和能
%% point: 动能定理
%% method: 无
%% score: 12
%% degree: 450
%% duplicateId: 0
%% body:
风洞是研究空气动力学的实验设备。如图，将刚性杆水平固定在风洞内距地面高度 $H = 3.2\Um$ 处，杆上套一质量 $m = 3\Ukg$，可沿杆滑动的小球。将小球所受的风力调节为 $F = 15\UN$，方向水平向左。小球以速度 $v_{0} = 8\Ums$ 向右离开杆端，假设小球所受风力不变，取 $g = 10\Umsq$。求：
\begin{enumerate}
	\item
	小球落地所需时间和离开杆端的水平距离；
	\item
	小球落地时的动能。
	\item
	小球离开杆端后经过多少时间动能为 $78\UJ$？
\end{enumerate}
\onepicture[w=4.8cm, align=right]{\includesvg{figs/31}}
%% answer: （1）$t=0.8\Us$，$s=4.8\Um$（2）$E_{kt}=120\UJ$（3）$t_{1}=0.4\Us$，$t_{2}=0.24\Us$
\jdanswer{
\begin{enumerate}
	\item
	$t=0.8\Us$，$s=4.8\Um$

	\item
	$E_{kt}=120\UJ$

	\item
	$t_{1}=0.4\Us$，$t_{2}=0.24\Us$

\end{enumerate}
}
%% memo:
\memoanswer{
(1) 小球在竖直方向做自由落体运动，运动时间为

$t = \sqrt{\frac{2H}{g}} = 0.8\Us$，

小球在水平方向做匀减速运动，加速度大小为

$a = \frac{F}{m} = 5\Umsq$，

水平位移为

$s = v_{0}t - \frac{1}{2}at^{2} = 4.8\Um$。

(2) 由动能定理得

$E_{kt} - E_{k0} = mgH - Fs$，解得 $E_{kt} = 120\UJ$。

(3) 小球离开杆后经过时间 $t$ 的水平位移为

$s = v_{0}t - \frac{1}{2}at^{2}$，

由动能定理得

$E_{k} - \frac{1}{2}mv_{0}^{2} = mg\cdot\frac{1}{2}gt^{2} - Fs$，

以 $E_{k} = 78\UJ$ 和 $v_{0} = 8\Ums$ 代入得

$125t^{2} - 80t + 12 = 0$，

解得 $t_{1} = 0.4\Us$，$t_{2} = 0.24\Us$。
}

\item
%% number: 32
%% paperName: 2016年上海高考第 32 题 14 分
%% typeId: 6
%% type: 计算题
%% chapter: 电场
%% point: 电场中的图象
%% method: 无
%% score: 14
%% degree: 350
%% duplicateId: 0
%% body:
如图~\ref{2016上海32a}，长度 $L = 0.8\Um$ 的光滑杆左端固定一带正电的点电荷 A，其电荷量 $Q = 1.8\times10^{-7}\UC$；一质量 $m = 0.02\Ukg$，带电量为 $q$ 的小球 B 套在杆上。将杆沿水平方向固定于某非均匀外电场中，以杆左端为原点，沿杆向右为 $x$ 轴正方向建立坐标系。点电荷 A 对小球 B 的作用力随 B 位置 $x$ 的变化关系如图~\ref{2016上海32b} 中曲线 I 所示，小球 B 所受水平方向的合力随 B 位置 $x$ 的变化关系如图~\ref{2016上海32b} 中曲线 II 所示，其中曲线 II 在 $0.16 \leq x \leq 0.20$ 和 $x \geq 0.40$ 范围可近似看作直线。求：（静电力常量 $k = 9\times10^{9}\UNmqCq$）
\begin{enumerate}
	\item
	小球 B 所带电量 $q$；
	\item
	非均匀外电场在 $x = 0.3\Um$ 处沿细杆方向的电场强度大小 $E$；
	\item
	在合电场中，$x = 0.4\Um$ 与 $x = 0.6\Um$ 之间的电势差 $U$。
	\item
	已知小球在 $x = 0.2\Um$ 处获得 $v = 0.4\Ums$ 的初速度时，最远可以运动到 $x = 0.4\Um$ 处。若小球在 $x = 0.16\Um$ 处受到方向向右，大小为 $0.04\UN$ 的恒力作用后，由静止开始运动，为使小球能离开细杆，恒力作用的最小距离 $s$ 是多少？
\end{enumerate}
\twopicture[labelA=2016上海32a, labelB=2016上海32b, numA=（a）, numB=（b）, wA=5.4cm, wB=6.2cm, align=right]
{\includesvg{figs/32a}}
{\includesvg{figs/32b}}
%% answer: （1）$q=1\times10^{-6}\UC$（2）$E=3\times10^{4}\UNC$，方向水平向左（3）$U=800\UV$（4）$s=0.065\Um$
\jdanswer{
\begin{enumerate}
	\item
	$q=1\times10^{-6}\UC$

	\item
	$E=3\times10^{4}\UNC$，方向水平向左

	\item
	$U=800\UV$

	\item
	$s=0.065\Um$

\end{enumerate}
}
%% memo:
\memoanswer{
(1) 由图可知，当 $x = 0.3\Um$ 时，

$F_{1} = k\frac{qQ}{x^{2}} = 0.018\UN$，解得 $q = \frac{F_{1}x^{2}}{kQ} = 1\times10^{-6}\UC$。

(2) 设在 $x = 0.3\Um$ 处点电荷与小球间作用力为 $F_{2}$，由力的合成知

$F_{\text{合}} = F_{2} + qE$，

解得

$E = \frac{F_{\text{合}} - F_{2}}{q} = \frac{-0.012 - 0.018}{1\times10^{-6}}\UNC = -3\times10^{4}\UNC$，

电场在 $x = 0.3\Um$ 处沿细杆方向的电场强度大小为 $3\times10^{4}\UNC$，方向水平向左。

(3) 由图象可知在 $x = 0.4\Um$ 与 $x = 0.6\Um$ 之间合力做功大小为

$W_{\text{合}} = 0.004\times0.2\UJ = 8\times10^{-4}\UJ$，

又 $qU = W_{\text{合}}$，

解得 $U = \frac{W_{\text{合}}}{q} = 800\UV$。

(4) 由图可知小球从 $x=0.16\Um$ 到 $x=0.2\Um$ 处，电场力做功的大小为

$W_{1}=\frac{0.03\times0.04}{2}\UJ=6\times10^{-4}\UJ$，

小球从 $x=0.2\Um$ 到 $x=0.4\Um$ 处，电场力做功的大小为

$W_{2}=-\frac{1}{2}mv^{2}=-1.6\times10^{-3}\UJ$，

由图可知小球从 $x=0.4\Um$ 到 $x=0.8\Um$ 处，电场力做功的大小为

$W_{3}=-0.004\times0.4\UJ=-1.6\times10^{-3}\UJ$，

由动能定理得

$W_{1}+W_{2}+W_{3}+F_{\text{外}}s=0$，

解得 $s=-\frac{W_{1}+W_{2}+W_{3}}{F_{\text{外}}}=0.065\Um$。
}

\item
%% number: 33
%% paperName: 2016年上海高考第 33 题 14 分
%% typeId: 6
%% type: 计算题
%% chapter: 电磁感应
%% point: 电磁感应中的匀变速
%% method: 无
%% score: 14
%% degree: 250
%% duplicateId: 0
%% body:
如图，一关于 $y$ 轴对称的导体轨道位于水平面内，磁感应强度为 $B$ 的匀强磁场与平面垂直。一足够长，质量为 $m$ 的直导体棒沿 $x$ 轴方向置于轨道上，在外力 $F$ 作用下从原点由静止开始沿 $y$ 轴正方向做加速度为 $a$ 的匀加速直线运动，运动时棒与 $x$ 轴始终平行。棒单位长度的电阻为 $\rho$，与电阻不计的轨道接触良好，运动中产生的热功率随棒位置的变化规律为 $P=ky^{3/2}$（SI）。求：
\begin{enumerate}
	\item
	导体轨道的轨道方程 $y=f(x)$；
	\item
	棒在运动过程中受到的安培力 $F_{m}$ 随 $y$ 的变化关系；
	\item
	棒从 $y=0$ 运动到 $y=L$ 过程中外力 $F$ 的功。
\end{enumerate}
\onepicture[w=4.6cm, align=right]{\includesvg{figs/33}}
%% answer: （1）$y=\left(\frac{4aB^{2}}{k\rho}\right)^{2}x^{2}$，轨道形状为抛物线（2）$F_{m}=\frac{k}{\sqrt{2a}}y$（3）$W=\frac{k}{2\sqrt{2a}}L^{2}+maL$
\jdanswer{
\begin{enumerate}
	\item
	$y=\left(\frac{4aB^{2}}{k\rho}\right)^{2}x^{2}$，轨道形状为抛物线

	\item
	$F_{m}=\frac{k}{\sqrt{2a}}y$

	\item
	$W=\frac{k}{2\sqrt{2a}}L^{2}+maL$

\end{enumerate}
}
%% memo:
\memoanswer{
(1) 设棒运动到某一位置时与轨道接触点的坐标为 $(\pm x, y)$，棒受到的安培力大小为 $F = \frac{B^{2}l^{2}v}{R}$，

安培力的功率为

$P = \frac{4B^{2}x^{2}v^{2}}{R} = ky^{\frac{3}{2}}$，

棒做匀加速运动时，由运动学公式得

$v^{2}=2ay$，

回路的总电阻为

$R = 2\rho x$，

联立解得 $y = \left(\frac{4aB^{2}}{k\rho}\right)^{2}x^{2}$，

故轨道的形状为抛物线。

(2) 棒受到的安培力大小为

$F_{m}=\frac{4B^{2}x^{2}}{R}v=\frac{2B^{2}x}{\rho}\sqrt{2ay}$，

以轨道方程代入得

$F_{m}=\frac{k}{\sqrt{2a}}y$。

(3) 由动能定理得

$W = W_{m} + \frac{1}{2}mv^{2}$，

安培力做的功为

$W_{m}=\frac{k}{2\sqrt{2a}}L^{2}$，

棒在 $y=L$ 处动能为

$\frac{1}{2}mv^{2}=maL$，

故外力做的功为

$W=\frac{k}{2\sqrt{2a}}L^{2}+maL$。
}

\end{enumerate}

\end{document}
